\documentclass{article}
\usepackage{amsmath,amssymb,amsfonts}
\usepackage{geometry}
\usepackage{hyperref}
\usepackage[ruled,vlined]{algorithm2e}
\usepackage{booktabs}
\usepackage{tabularx}
\usepackage{xcolor}
\geometry{a4paper, margin=1in}

\newcommand{\code}[1]{\texttt{\small #1}}
\newcommand{\textftol}{\tau_{\mathrm{ftol}}}
\newcommand{\textstag}{\tau_{\mathrm{stag}}}
\newcommand{\textxtol}{\tau_{\mathrm{xtol}}}
\newcommand{\ftolinner}[1]{\tau^{\,\mathrm{inner}}_{\mathrm{ftol},#1}}

\title{Algorithmic Structure of the Stabilized Porous Navier--Stokes Solver}
\author{Antigravity AI \and Project Authors}
\date{\today}

\begin{document}
\maketitle

\begin{abstract}
This document specifies the complete algorithmic pipeline that drives a single
invocation of \code{run\_simulation} in the stabilized porous Navier--Stokes
solver. The pipeline is presented as a cascade of seven nested
algorithms---one outermost orchestration box, three for the inner Newton
machinery, one for the Picard globalization fallback, one for the staggered
Orthogonal Sub-Grid Scale (OSGS) outer loop, and one for the Method of
Manufactured Solutions (MMS) plateau verifier. For each algorithm we provide
both rigorous pseudocode and pedagogical exposition: what the algorithm does,
why each branch exists, how its numerical parameters influence cost and
robustness, and how it couples to its neighbours in the cascade. A separate
section documents three adaptations that live exclusively in the MMS sweep
harness and do \emph{not} participate in single-shot \code{run\_simulation}
invocations.
\end{abstract}

\tableofcontents

% ============================================================
\section{Overview and Algorithmic Challenges}
\label{sec:overview}

The continuous problem is the stationary, incompressible, generalized porous
Navier--Stokes system on a domain $\Omega \subset \mathbb{R}^d$,
\begin{equation}
    \label{eq:strong}
    \begin{aligned}
    \boldsymbol{u} \cdot \nabla \boldsymbol{u}
        - \nabla \cdot (2 \nu \nabla^{\mathrm{s}} \boldsymbol{u})
        + \sigma(\alpha, \boldsymbol{u}) \boldsymbol{u}
        + \nabla p
        &= \boldsymbol{f}, \\
    \nabla \cdot \boldsymbol{u} &= 0,
    \end{aligned}
\end{equation}
discretised by an equal-order Lagrangian pair $(\boldsymbol{V}_h, Q_h)$ of
degree $(k_v, k_p) = (k, k)$ and stabilised by a Variational Multiscale (VMS)
formulation. The drag tensor $\sigma(\alpha, \boldsymbol{u})
= a(\alpha) + b(\alpha) |\boldsymbol{u}|$ couples the Darcy--Forchheimer regime
to the convective momentum balance; the porosity field $\alpha$ varies in
space, and the high-$\mathrm{Re}$, high-$\mathrm{Da}$ regimes the solver must
handle are exactly those in which the system becomes simultaneously
convection-dominated, reaction-dominated, and stiff.

Two algorithmic challenges follow.

\paragraph{Nonlinearity from convection and Forchheimer drag.}  Newton's
method converges quadratically once the iterate enters its quadratic basin,
but locating that basin from an arbitrary initial guess is non-trivial when
both $\boldsymbol{u}\cdot\nabla\boldsymbol{u}$ and the $|\boldsymbol{u}|$
factor in $\sigma$ amplify any error in $\boldsymbol{u}$. Outside the basin,
a raw Newton step typically over-corrects, and the merit function
$\Phi(\boldsymbol{x}) = \tfrac{1}{2}\|\boldsymbol{b}(\boldsymbol{x})\|^2$
grows. The solver must detect this divergence and respond---either by
shortening the step (line search), by switching the iteration map to one with
a wider basin of attraction (Picard), or by aborting cleanly.

\paragraph{Implicit OSGS projection.}  In the Orthogonal Sub-Grid Scale
formulation the sub-grid component lives in the orthogonal complement
$\mathcal{X}_{h0}^{\perp}$ of the finite element space; in practice the
projection $\boldsymbol{\pi}_h \in \mathcal{X}_h$ of the strong residual is
computed, and the sub-grid scale is identified with $(I-\Pi_h)\mathcal{R}$.
The projection itself depends on the primary fields $U_h = (\boldsymbol{u}_h,
p_h)$, so a fully monolithic Jacobian would embed the
Fr\'echet derivative of an $L^2$ projection operator. Rather than carrying
this dense block, the solver freezes $\boldsymbol{\pi}_h$, solves the
nonlinear primary system, then re-projects, in a staggered (block-coordinate)
fixed point.

\paragraph{Architecture of the response.}  The architecture that emerges is a
two-stage cascade. \emph{Stage~I} drives the algebraic system to a converged
ASGS state, regardless of the final stabilization method requested. It
composes an exact-Newton attempt, a Picard fallback, and a final Newton
polish (Algorithms~\ref{alg:A} and~\ref{alg:B}). \emph{Stage~II} either
reports the ASGS state directly or, when the user has selected OSGS, runs
a staggered outer loop on $(U_h, \boldsymbol{\pi}_h)$ initialised at the
Stage~I solution (Algorithm~\ref{alg:C}). An optional MMS verifier
(Algorithm~\ref{alg:D}) extends each Stage~II loop with a plateau test that
distinguishes \emph{algebraic} convergence (small nonlinear residual) from
\emph{discretisation} convergence (small error against an exact solution).
The outermost orchestration (Algorithm~\ref{alg:O}) wires all of these
together. The reader who wants a single mental model should hold:

\begin{quote}
\itshape
ASGS is always run first; OSGS, when requested, refines the ASGS state by
a staggered projection loop; MMS verification, when enabled, runs as a
plateau gate sitting at the top of whichever Stage~II loop is active.
\end{quote}

% ============================================================
\section{Parameter Taxonomy and Code Equivalence}
\label{sec:params}

\subsection{Numerical parameters}

Table~\ref{tab:params} enumerates every numerical parameter consumed by the
algorithms below, grouped by role. Parameters marked \emph{harness-only} live
in the configuration schema but are read \emph{only} by the MMS sweep harness
(\code{test/extended/ManufacturedSolutions/run\_test.jl}); see
Section~\ref{sec:harness}. All remaining parameters participate in a single
\code{run\_simulation} call.

\begin{table}[h]
\centering
\caption{Numerical parameters of the nonlinear solver pipeline. The
``Effect'' column summarises the dominant influence of each parameter on
robustness or cost. Each parameter participates in every \code{run\_simulation}
call except those flagged \emph{harness-only}; see Section~\ref{sec:harness}.}
\label{tab:params}
\footnotesize
\begin{tabularx}{\textwidth}{@{} l l X @{}}
\toprule
\textbf{Symbol} & \textbf{Field in code} & \textbf{Role / effect} \\
\midrule
\multicolumn{3}{l}{\emph{Inner Newton iteration controls}} \\
$N_{\mathrm{N}}$         & \code{sol.newton\_iterations}           & Cap on inner Newton iterations. \\
$N_{\mathrm{P}}$         & \code{sol.picard\_iterations}           & Cap on Picard fallback iterations. \\
$N_{\mathrm{inc}}$       & \code{sol.max\_increases}               & Consecutive merit-function expansions tolerated before declaring divergence. Smaller is stricter; larger is more permissive at the cost of slower failure detection. \\
$\textftol$              & \code{sol.ftol}                          & Terminal residual tolerance in $\ell^\infty$; the dominant convergence criterion. \\
$\textxtol$              & \code{sol.xtol}                          & Lower bound on the accepted step norm $\|\alpha\delta\boldsymbol{x}\|$ before declaring coordinate stagnation. \\
$\textstag$              & \code{sol.stagnation\_noise\_floor}      & Noise-floor residual; a non-converged exit below this may still be classified as success (Sec.~\ref{sec:falseEquiv}). \\
$\alpha_{\min}$          & \code{sol.linesearch\_alpha\_min}        & Smallest line-search step before giving up. Smaller means more aggressive descent, but more failed Newton iterations per outer call. \\
$c_1$                    & \code{sol.armijo\_c1}                    & Armijo descent constant. Smaller means easier step acceptance and a weaker guarantee. \\
$\gamma$                 & \code{sol.linesearch\_contraction\_factor} & Backtracking factor: $\alpha \leftarrow \gamma\alpha$. Smaller tries fewer step sizes but explores more coarsely. \\
$N_{\mathrm{ls}}$        & \code{sol.max\_linesearch\_iterations}   & Bound on line-search inner iterations. \\
$F_{\mathrm{div}}$       & \code{sol.divergence\_merit\_factor}     & Threshold by which the merit must expand to count as one ``increase''. Smaller is a stricter divergence detector. \\
\code{freeze\_cusp} & \code{sol.freeze\_jacobian\_cusp}   & Boolean: skip Jacobian recomputation at Forchheimer cusps to avoid singular updates. \\
\midrule
\multicolumn{3}{l}{\emph{OSGS staggered outer loop}} \\
$N_{\mathrm{OSGS}}$      & \code{stab.osgs\_iterations}             & Cap on outer (staggered) iterations. \\
$N_{\mathrm{inner}}$     & \code{stab.osgs\_inner\_newton\_iters}   & Cap on inner Newton iterations per outer step. Smaller makes each sweep cheap and forces more outer iterations; larger means a deeper inner solve and fewer outer iterations. \\
$\tau_{\mathrm{OSGS}}$   & \code{stab.osgs\_tolerance}              & Tolerance on the state drift $d_{\mathrm{state}}^m$. \\
$\tau_{\mathrm{proj}}$   & \code{stab.osgs\_projection\_tolerance}  & Tolerance on the mass-weighted projection drift $d_{\pi}^m$. \\
$m_{\mathrm{w}}$         & \code{stab.osgs\_warmup\_iterations}     & Outer iterations during which the inner Newton uses a relaxed tolerance, to avoid polishing iterates against a stale projection. \\
$\tau_{\mathrm{w}}$      & \code{stab.osgs\_warmup\_tolerance}      & Relaxed inner tolerance during warmup; the effective inner tolerance is $\max(\textftol,\tau_{\mathrm{w}})$. \\
$\mathsf{Mode}_{\mathrm{drift}}$ & \code{stab.osgs\_state\_drift\_scale} & Drift metric: \code{Linf} (raw $\ell^\infty$ on DOFs) or \code{L2\_mass} (per-block continuous $L^2(\Omega)$). See Sec.~\ref{sec:drift-modes}. \\
$\mathsf{Mode}_{\mathrm{stop}}$  & \code{stab.osgs\_stopping\_mode}      & Convergence predicate: \code{state\_drift}, \code{projection\_drift}, or \code{both} (conjunction; the default). \\
$m$, $\beta$             & \code{sol.accelerator.\{m,relaxation\_factor\}} & Anderson history depth and damping. The accelerator is enabled by setting \code{sol.accelerator.type} to \code{"Anderson"}. \\
\midrule
\multicolumn{3}{l}{\emph{MMS plateau verification}} \\
$N_{\mathrm{plat}}$      & \code{mms.require\_consecutive\_passes}  & Consecutive passes of $r_{\max} \le \tau_{\mathrm{err}}$ required before declaring a plateau. \\
$\tau_{\mathrm{err}}$    & \code{mms.tau\_err}                       & Relative-change tolerance against which $r_{\max}$ is tested. \\
$\varepsilon_{u,L^2}, \varepsilon_{u,H^1}, \varepsilon_{p,L^2}$ & \code{mms.eps\_u\_l2,\_u\_h1,\_p\_l2} & Per-norm floors guarding the denominator in $r_{\max}$. \\
$K_{\mathrm{extra}}$     & \code{mms.max\_extra\_cycles}             & Cap on extension cycles (ASGS path) or staggered iterations added to the OSGS budget (OSGS path) for plateau verification. \\
\midrule
\multicolumn{3}{l}{\emph{Harness-only adaptations} (Section~\ref{sec:harness})} \\
$C_{\mathrm{ceil}}$, $C_{\mathrm{sf}}$ & \code{sol.dynamic\_ftol\_ceiling}, \code{\_spatial\_safety\_factor} & Clamps for the adaptive $\textftol$ rule of Sec.~\ref{sec:harness-ftol}. \\
$n^{\mathrm{base}}, n^{\min}, n^{\mathrm{sf}}$ & \code{sol.condition\_noise\_floor\_*} & Adaptive noise-floor controls (Sec.~\ref{sec:harness-noise}). \\
$\mathrm{Re}^{\mathrm{th}}, N_{\mathrm{P}}^{\mathrm{Re}}$ & \code{sol.dynamic\_picard\_re\_*} & Per-case Picard widening trigger and budget for high Reynolds. \\
$\mathrm{Da}^{\mathrm{th}}, N_{\mathrm{P}}^{\mathrm{Da}}$ & \code{sol.dynamic\_picard\_da\_*} & Per-case Picard widening trigger and budget for high Damk\"ohler. \\
\bottomrule
\end{tabularx}
\end{table}

\subsection{Algorithm--to--code correspondence}

Table~\ref{tab:functions} pins each algorithm presented below to its Julia
entry point.

\begin{table}[h]
\centering
\caption{Algorithm $\leftrightarrow$ code signature mapping.}
\label{tab:functions}
\small
\begin{tabularx}{\textwidth}{@{} l X @{}}
\toprule
\textbf{Algorithm} & \textbf{Code entry point} \\
\midrule
$\mathrm{SimulationOrchestration}$ (Alg.~\ref{alg:O})   & \code{run\_simulation(path)} $\to$ \code{solve\_system(setup, form, solvers, cfg, x0)} (\code{src/run\_simulation.jl}, \code{src/solvers/porous\_solver.jl}). \\
$\mathrm{ExactNewtonPipeline}$ (Alg.~\ref{alg:A})       & \code{Gridap.Algebra.solve!(x, FESolver(SafeNewtonSolver(\dots)), op)} dispatching into \code{\_safe\_solve\_inner!} (\code{src/solvers/nonlinear.jl}). The same routine is reused with an \code{ExactNewtonMode()} Jacobian or a Picard Jacobian; the safeguard logic does not change. \\
$\mathrm{SolveLinearSystem}$ (Alg.~\ref{alg:A1})        & \code{eval\_linear\_system\_resolution!(dx, A, b, solver, cache)}. \\
$\mathrm{ArmijoLinesearch}$ (Alg.~\ref{alg:A2})         & \code{eval\_armijo\_linesearch\_pass!(x, x\_old, b, dx, op, dir\_deriv, solver, w, state)}. \\
$\mathrm{EvaluateSafeguards}$ (Alg.~\ref{alg:A3})       & \code{eval\_safeguard\_termination\_bounds!(solver, state)}. \\
$\mathrm{PicardFallback}$ (Alg.~\ref{alg:B})            & The Newton$\to$Picard$\to$Newton cascade in \code{solve\_system}, lines 169--246 of \code{src/solvers/porous\_solver.jl}. \\
$\mathrm{OSGSFractionalRelaxation}$ (Alg.~\ref{alg:C})  & The \code{for osgs\_iter in 1:eff\_osgs\_iters} block in \code{solve\_system}, lines 340--615. \\
$\mathrm{VerifyMMSPlateau}$ (Alg.~\ref{alg:D})          & The extended-cycle loop in \code{solve\_system} (ASGS entry, lines 258--337) and the in-loop plateau gate (OSGS entry, lines 546--587). \\
\bottomrule
\end{tabularx}
\end{table}

% ============================================================
\section{The Outermost Orchestration: \texorpdfstring{Algorithm~\ref{alg:O}}{Algorithm O}}
\label{sec:orchestration}

The orchestration drives one simulation end-to-end. Its central organising
principle is the two-stage cascade described in
Section~\ref{sec:overview}: an unconditional ASGS algebraic
initialisation followed by a method-dependent dispatch. Notice that
Stage~I runs even when the user has requested OSGS, because the staggered
projection loop is much better conditioned when initialised at a converged
ASGS state than at a zero initial guess.

\begin{algorithm}[H]
\caption{$\mathrm{SimulationOrchestration}$: end-to-end driver.}
\label{alg:O}
\SetAlgoLined
\KwIn{Configuration path $P$.}
\KwOut{$(u_h, p_h, \boldsymbol{\pi}_h)$, success flag $S$, total inner-iteration count.}
\vspace{0.2cm}\hrule\vspace{0.2cm}

$\Theta \gets \mathrm{load\_config\_with\_defaults}(P)$ \tcp*{validate; reject unknown keys}
$M \gets \mathrm{build\_default\_mesh}(\Theta)$\;
$(X, Y, k_v, k_p) \gets \mathrm{build\_fe\_spaces}(M, \Theta)$ \tcp*{Dirichlet-constrained trial/test pairs}
$(V_{\mathrm{free}}, Q_{\mathrm{free}}) \gets$ unconstrained mirrors of $(X,Y)$ \tcp*{projection spaces}
$F \gets \mathrm{build\_formulation}(\Theta)$ \tcp*{viscous + reaction + $\tau$ + projection policy}
$T \gets \mathrm{FETopology}(X,Y,M,\Omega,\mathrm{d}\Omega,V_{\mathrm{free}},Q_{\mathrm{free}}, h, \boldsymbol{f}, \alpha)$\;
$\boldsymbol{x}^{0} \gets \mathrm{interpolate\_everywhere}(0, X)$\;

\vspace{0.15cm}
\tcc{\textbf{Stage I.} Algebraic initialisation in ASGS mode ($\boldsymbol{\pi}_h \equiv 0$).}
$(S, U_h^{\mathrm{ASGS}}) \gets \mathrm{PicardFallback}(\boldsymbol{x}^0,\ \boldsymbol{x}^0_{\mathrm{backup}})$ \tcp*{Alg.~\ref{alg:B}}
\If{\textnormal{\textbf{not }} $S$}{
    \Return failure (no fields written)\;
}

\vspace{0.15cm}
\tcc{\textbf{Stage II.} Method dispatch.}
\If{$\Theta.\mathrm{method} = \mathrm{ASGS}$}{
    \If{$\Theta.\mathrm{mms}.\mathrm{enabled}$}{
        run $K_{\mathrm{extra}}$ extension cycles, calling $\mathrm{VerifyMMSPlateau}$ each cycle (Alg.~\ref{alg:D}, ASGS entry)\;
    }
    \Return $(U_h^{\mathrm{ASGS}}, \boldsymbol{\pi}_h \equiv 0, \dots)$\;
}
\If{$\Theta.\mathrm{method} = \mathrm{OSGS}$}{
    pre-assemble and LU-factor $M_u = (V_{\mathrm{free}}, U_{\mathrm{free}})_{L^2}$, $M_p = (Q_{\mathrm{free}}, P_{\mathrm{free}})_{L^2}$ once\;
    $\boldsymbol{\pi}_h^{0} \gets L^2$-projection of $\mathcal{R}(U_h^{\mathrm{ASGS}})$ onto $(V_{\mathrm{free}}, Q_{\mathrm{free}})$\;
    run $\mathrm{OSGSFractionalRelaxation}(U_h^{\mathrm{ASGS}}, \boldsymbol{\pi}_h^{0})$ (Alg.~\ref{alg:C}) \tcp*{embeds Alg.~\ref{alg:D} as plateau gate}
}

\vspace{0.15cm}
$\mathrm{export\_results}(\mathrm{VTK} \mid \mathrm{HDF5})$\;
\end{algorithm}

\paragraph{Step by step.}  Lines~1--7 load and validate the configuration,
build the mesh, and assemble the finite-element pair. The unconstrained
spaces $V_{\mathrm{free}}, Q_{\mathrm{free}}$ are required \emph{only} by the
OSGS projection step (Section~\ref{sec:osgs}); they carry no Dirichlet data,
because projecting onto a Dirichlet-constrained space introduces an $O(1)$
boundary residual that destroys the expected $O(h^{k+1})$ MMS rate. The
\code{FETopology} bundle simply ties together the geometry, the test/trial
spaces, the per-element mesh size $h$, and the loading data
$(\boldsymbol{f}, \alpha)$.

Stage~I (line~8) is the workhorse: it calls $\mathrm{PicardFallback}$
(Algorithm~\ref{alg:B}), which composes an exact-Newton attempt, a
Picard globalisation pass, and a final Newton polish, all with
$\boldsymbol{\pi}_h \equiv 0$. The output $U_h^{\mathrm{ASGS}}$ is the
algebraic root of the ASGS system; when Stage~II is OSGS, the staggered loop
will descend from this state with an already-converged primary field.

The Stage~II dispatch (lines~11--19) is intentionally simple: either accept
the ASGS state, optionally extended by an MMS plateau check, or hand off to
the OSGS outer loop. The mass matrices $M_u, M_p$ are pre-assembled and
LU-factored exactly once before the OSGS loop begins (line~16). Because the
projection target spaces $V_{\mathrm{free}}, Q_{\mathrm{free}}$ never change
during the loop, these factorisations are reused at every outer iteration,
which converts every $L^2$ projection into a constant-cost forward/back
substitution.

\paragraph{Failure propagation.}  If Stage~I returns $S = \mathrm{False}$,
the orchestrator returns immediately without writing output (lines~9--10).
The diagnostics cache carries the algebraic exit reason so the calling
script can distinguish a divergent run from a converged one.

\paragraph{Anderson and MMS coupling.}  Two side-effects of $\Theta$ deserve
mention. First, if the configured accelerator type is \code{"Anderson"},
the OSGS loop constructs a pair of block-wise Anderson accelerators
(Section~\ref{sec:anderson}); the orchestrator itself does not allocate
them, but their presence changes the iterate produced at each outer step.
Second, when $\Theta.\mathrm{mms}.\mathrm{enabled}$, the OSGS outer-iteration
budget $N_{\mathrm{OSGS}}$ is silently inflated by $K_{\mathrm{extra}}$ so
that the plateau verifier has room to run after the staggered fixed point
itself has converged. Without this inflation, the verifier would chronically
exhaust the budget on the very first extension cycle.

% ============================================================
\section{The Inner Newton Machinery}
\label{sec:newton}

\subsection{The merit function and its diagonal preconditioning}
\label{sec:merit}

Let $\boldsymbol{b}(\boldsymbol{x}) \in \mathbb{R}^N$ be the assembled residual
of the (possibly stabilised) discrete weak form, and $\mathcal{J}(\boldsymbol{x})
= \partial \boldsymbol{b}/\partial \boldsymbol{x}$ its Jacobian. The merit
function used by the line search is the diagonally equilibrated quadratic
\begin{equation}
    \label{eq:merit}
    \Phi_{\boldsymbol{w}}(\boldsymbol{x})
        = \tfrac{1}{2} \sum_{k=1}^{N}
            \left( \frac{b_k(\boldsymbol{x})}{w_k} \right)^{2},
    \qquad
    w_k = \max\!\bigl(|\mathcal{J}_{kk}|,\ \varepsilon_{\mathrm{mach}}
            \|\mathrm{diag}(\mathcal{J})\|_\infty,\ \varepsilon_{\mathrm{mach}} \bigr).
\end{equation}
The weights $w_k$ are recomputed at every Newton iteration from the diagonal
of the current Jacobian.

\paragraph{Why equilibrate?}  The discrete residual mixes momentum-balance
rows---measured in physical units of force per unit volume, and possibly
amplified by a large Reynolds or Damk\"ohler number---with mass-balance rows,
whose natural magnitude is dimensionally different. A raw unweighted merit
$\tfrac{1}{2}\|\boldsymbol{b}\|^2$ would let the momentum rows dominate
entirely, so the line search would be quietly indifferent to mass-conservation
errors. Equilibrating by $w_k$ rescales every row to its own natural unit, and
the line search now respects both balances simultaneously. The lower bound
$\varepsilon_{\mathrm{mach}}\|\mathrm{diag}(\mathcal{J})\|_\infty$ prevents a
genuinely small diagonal entry (e.g.\ a pressure row when the system is
saddle-point-like) from triggering division by zero.

\paragraph{An algebraic identity that simplifies the line search.}  For a
Newton step $\delta\boldsymbol{x}$ satisfying $\mathcal{J}\delta\boldsymbol{x}
= \boldsymbol{b}$, the directional derivative of $\Phi_{\boldsymbol{w}}$
along $-\delta\boldsymbol{x}$ at $\boldsymbol{x}$ is
\begin{equation}
    \label{eq:dir-deriv}
    D \;=\; -\sum_k w_k^{-2} b_k(\mathcal{J}\delta\boldsymbol{x})_k
        \;=\; -\sum_k w_k^{-2} b_k^2
        \;=\; -2\,\Phi_{\boldsymbol{w}}(\boldsymbol{x}).
\end{equation}
The code exploits this identity directly: \code{dir\_deriv = -2.0 * state.phi\_x}
at \code{nonlinear.jl:216}. The implication for Algorithm~\ref{alg:A2} is
that the Armijo predicate $\hat\Phi \le \Phi + c_1\alpha D$ becomes
$\hat\Phi \le \Phi(1 - 2 c_1 \alpha)$, which never requires an extra
gradient evaluation. The identity holds \emph{only} because the search
direction is the exact Newton step; a Picard step in general violates it,
which is one reason the Picard fallback in Algorithm~\ref{alg:B} runs with
its own (cruder) line-search settings.

\subsection{Algorithm A: the safeguarded Newton outer loop}

\begin{algorithm}[H]
\caption{$\mathrm{ExactNewtonPipeline}$: safeguarded Newton iteration with merit-based line search.}
\label{alg:A}
\SetAlgoLined
\textbf{Parameters:} $N_{\mathrm{N}}$, $\textftol$, $\textxtol$, $\textstag$, $c_1$, $\gamma$, $\alpha_{\min}$, $N_{\mathrm{ls}}$, $F_{\mathrm{div}}$, $N_{\mathrm{inc}}$.\\
\KwIn{Initial fields $\boldsymbol{x}^0 \in \mathcal{X}_h$, nonlinear operator $\mathrm{op}$ (encoding either an exact-Newton or a Picard Jacobian).}
\KwOut{$(\boldsymbol{x}^{\star}, S \in \{\mathrm{True},\mathrm{False}\}, I, \|\boldsymbol{b}^{\star}\|_\infty, r)$, where $r$ is the exit reason (Table~\ref{tab:reasons}).}
\vspace{0.2cm}\hrule\vspace{0.2cm}

$\boldsymbol{b} \gets \boldsymbol{b}(\boldsymbol{x}^0)$;\quad
\If{$\|\boldsymbol{b}\|_\infty \le \textftol$}{
    \Return $(\boldsymbol{x}^0,\ \mathrm{True},\ 0,\ \|\boldsymbol{b}\|_\infty,\ \mathsf{InitialFtol})$\;
}
$\boldsymbol{x}_{\mathrm{best}} \gets \boldsymbol{x}^0$;\quad
$b_{\mathrm{best}} \gets \|\boldsymbol{b}\|_\infty$;\quad
$I_{\mathrm{up}} \gets 0$;\quad
$r \gets \mathsf{IterationCap}$\;
\For{$i = 1, \dots, N_{\mathrm{N}}$}{
    $\mathcal{J} \gets \mathcal{J}(\boldsymbol{x})$;\quad
    update weights $\boldsymbol{w}$ via~\eqref{eq:merit};\quad
    $\Phi \gets \Phi_{\boldsymbol{w}}(\boldsymbol{x})$\;
    $(S_{\mathrm{LS}}^{\mathrm{lin}}, \delta\boldsymbol{x}) \gets \mathrm{SolveLinearSystem}(\mathcal{J}, \boldsymbol{b})$ \tcp*{Alg.~\ref{alg:A1}}
    \If{$S_{\mathrm{LS}}^{\mathrm{lin}} = \mathrm{Failed}$}{
        $\boldsymbol{x} \gets \boldsymbol{x}_{\mathrm{best}}$;\quad
        \Return $(\boldsymbol{x},\ \mathrm{False},\ i,\ b_{\mathrm{best}},\ \mathsf{LinearSolveNaN})$\;
    }
    set $D \gets -2\Phi$;\quad
    $(S_{\mathrm{LS}}, \alpha, L, \hat b, \hat \Phi)
        \gets \mathrm{ArmijoLinesearch}(\boldsymbol{x},\boldsymbol{x}_{\mathrm{old}},
              \boldsymbol{b},\delta\boldsymbol{x}, D, \boldsymbol{w}, \Phi)$ \tcp*{Alg.~\ref{alg:A2}}
    \If{$\hat{b} < b_{\mathrm{best}}$}{ $\boldsymbol{x}_{\mathrm{best}} \gets \boldsymbol{x}$;\quad $b_{\mathrm{best}} \gets \hat{b}$\; }
    $(r, S_{\mathrm{conv}}, S_{\mathrm{brk}}, I_{\mathrm{up}})
        \gets \mathrm{EvaluateSafeguards}(i, S_{\mathrm{LS}}, L, \hat b, \hat \Phi, \Phi, I_{\mathrm{up}})$ \tcp*{Alg.~\ref{alg:A3}}
    \If{$S_{\mathrm{brk}}$}{
        \If{$\neg S_{\mathrm{conv}}$ \textnormal{\textbf{and}} $r \notin \{\mathsf{Ftol}, \mathsf{NoiseFloor}, \mathsf{Xtol}\}$}{
            $\boldsymbol{x} \gets \boldsymbol{x}_{\mathrm{best}}$\;
        }
        \Return $(\boldsymbol{x},\ S_{\mathrm{conv}},\ i,\ \|\boldsymbol{b}^{\star}\|_\infty,\ r)$\;
    }
}
\Return $(\boldsymbol{x},\ \mathrm{False},\ N_{\mathrm{N}},\ \|\boldsymbol{b}\|_\infty,\ \mathsf{IterationCap})$\;
\end{algorithm}

\paragraph{Reading the box.}  The outer loop is a textbook safeguarded
Newton iteration: at each step compute the Jacobian, solve a linear system
for the Newton direction, search along it for a step that decreases the
merit, and ask the safeguard module whether to keep going. Three details
are worth highlighting because they are easy to miss.

\begin{itemize}
\item \emph{The safeguard's role is not just to detect convergence.} It also
detects \emph{progress without convergence} (residual is dropping; keep
going), \emph{stuck progress} (step is shrinking faster than the residual;
declare $\mathsf{Xtol}$ stagnation), and \emph{merit explosion} (the line
search succeeded but the new merit is far above the old; count an
``increase'' and abort if too many accumulate).

\item \emph{Best-state retention.}  Whenever the residual strictly improves,
$(\boldsymbol{x}_{\mathrm{best}}, b_{\mathrm{best}})$ is updated. On most
non-converged exits the code restores the iterate to this best-seen point,
so the caller never receives a state worse than what was already produced
during the iteration. Three exits intentionally \emph{do not} restore:
$\mathsf{Ftol}$ and $\mathsf{NoiseFloor}$ are already at or below the
target, and $\mathsf{Xtol}$ stagnation reflects a state the iterate could
not move beyond---restoring would discard the most current information.

\item \emph{Exact-Newton vs Picard reuse.}  The same loop is reused for the
Picard fallback in Algorithm~\ref{alg:B}; the only difference is the
Jacobian wired into \code{op}. The merit, line search, and safeguards apply
unchanged. This is what makes the cascade in Algorithm~\ref{alg:B} cheap to
maintain.
\end{itemize}

\subsection{Algorithm A.1: linear-system solve and NaN guard}

\begin{algorithm}[H]
\caption{$\mathrm{SolveLinearSystem}$.}
\label{alg:A1}
\SetAlgoLined
\textbf{Parameters:} (none).\\
\KwIn{Jacobian $\mathcal{J}$, residual $\boldsymbol{b}$.}
\KwOut{$(S_{\mathrm{Failed}} \in \{\mathrm{True},\mathrm{False}\}, \delta\boldsymbol{x})$.}
\vspace{0.2cm}\hrule\vspace{0.2cm}
\textbf{try} \{ Solve $\mathcal{J} \delta\boldsymbol{x} = \boldsymbol{b}$ via the configured linear solver \} \textbf{catch} \{ \Return $(\mathrm{True}, \delta\boldsymbol{x})$ \tcp*{linear-solver exception} \} \\
\If{any element of $\delta\boldsymbol{x}$ is NaN}{
    \Return $(\mathrm{True}, \delta\boldsymbol{x})$ \tcp*{catastrophic algebraic failure}
}
\Return $(\mathrm{False}, \delta\boldsymbol{x})$\;
\end{algorithm}

The linear solve is the cheapest place to detect a near-singular Jacobian:
either the factorisation aborts with an exception (the \code{try/catch} on
\code{nonlinear.jl:54--61}) or it returns a vector contaminated by NaNs.
Either way the outer loop must restore $\boldsymbol{x}_{\mathrm{best}}$ and
exit. Without this guard, a single NaN propagates through the residual
evaluation and silently destroys the iterate.

\subsection{Algorithm A.2: Armijo backtracking line search}

\begin{algorithm}[H]
\caption{$\mathrm{ArmijoLinesearch}$.}
\label{alg:A2}
\SetAlgoLined
\textbf{Parameters:} $\alpha_{\min}$, $c_1$, $\gamma$, $N_{\mathrm{ls}}$.\\
\KwIn{$\boldsymbol{x}$, $\boldsymbol{x}_{\mathrm{old}}$, $\boldsymbol{b}$, $\delta\boldsymbol{x}$, directional derivative $D = -2\Phi$, weights $\boldsymbol{w}$, current merit $\Phi$.}
\KwOut{$(S_{\mathrm{LS}} \in \{\mathrm{True},\mathrm{False}\}, \alpha, L, \|\hat{\boldsymbol{b}}\|_\infty, \hat\Phi)$.}
\vspace{0.2cm}\hrule\vspace{0.2cm}
$\alpha \gets 1$;\quad $S_{\mathrm{LS}} \gets \mathrm{False}$\;
\For{$j = 1, \dots, N_{\mathrm{ls}}$}{
    $\boldsymbol{x} \gets \boldsymbol{x}_{\mathrm{old}} - \alpha\,\delta\boldsymbol{x}$;\quad
    $\hat{\boldsymbol{b}} \gets \boldsymbol{b}(\boldsymbol{x})$;\quad
    $\hat{\Phi} \gets \Phi_{\boldsymbol{w}}(\hat{\boldsymbol{b}})$\;
    \If{$\hat\Phi$ is NaN}{
        $\alpha \gets \gamma \alpha$;\quad
        \textbf{continue} \tcp*{shrink and retry without testing Armijo}
    }
    \If{$\hat\Phi \le \Phi + c_1 \alpha D$}{
        $S_{\mathrm{LS}} \gets \mathrm{True}$;\quad \textbf{break}\;
    }
    \If{$\alpha \le \alpha_{\min}$}{
        \textbf{break} \tcp*{return with $S_{\mathrm{LS}} = \mathrm{False}$}
    }
    $\alpha \gets \gamma \alpha$\;
}
$L \gets \alpha \|\delta\boldsymbol{x}\|_\infty$\;
\Return $(S_{\mathrm{LS}}, \alpha, L, \|\hat{\boldsymbol{b}}\|_\infty, \hat\Phi)$\;
\end{algorithm}

\paragraph{What the parameters do.}
\begin{itemize}
\item $c_1$ controls how aggressive the descent must be. Setting $c_1$
small (say $10^{-4}$) accepts almost any step that lowers the merit at all;
setting $c_1$ close to $1$ demands a step that almost matches the full
linearised decrease. Newton implementations typically choose $c_1 \sim
10^{-4}$ because the Newton direction itself is steep enough.
\item $\gamma \in (0,1)$ is the contraction factor. A typical choice is
$\gamma = 1/2$; smaller $\gamma$ tries fewer step sizes but explores more
coarsely.
\item $\alpha_{\min}$ is the smallest step size accepted before declaring
the line search dead. Too small and the algorithm wastes iterations on
machine-epsilon steps; too large and it gives up on legitimately stiff
descent paths.
\item $N_{\mathrm{ls}}$ caps the inner iteration count and is in practice
the dominant cost knob: each iteration is a full residual evaluation.
\end{itemize}

\paragraph{NaN handling.}  When the residual evaluation at the trial point
produces a NaN merit (typically because the Forchheimer factor
$|\boldsymbol{u}|$ blew up under a too-aggressive step), the line search
shrinks $\alpha$ without testing Armijo. This is critical near the
quadratic-basin boundary: a Newton step pushing $\boldsymbol{u}$ into a
region where the model momentarily becomes non-evaluable should not be
recorded as a failed Armijo step; it should simply be retried at a smaller
size.

\subsection{Algorithm A.3: termination safeguards}

\begin{algorithm}[H]
\caption{$\mathrm{EvaluateSafeguards}$.}
\label{alg:A3}
\SetAlgoLined
\textbf{Parameters:} $\textftol$, $\textstag$, $\textxtol$, $F_{\mathrm{div}}$, $N_{\mathrm{N}}$, $N_{\mathrm{inc}}$.\\
\KwIn{Outer index $i$, line-search success $S_{\mathrm{LS}}$, step norm $L$, post-step residual $\|\hat{\boldsymbol{b}}\|_\infty$, post-step merit $\hat\Phi$, base merit $\Phi$, increase counter $I_{\mathrm{up}}$.}
\KwOut{$(r, S_{\mathrm{conv}}, S_{\mathrm{brk}}, I_{\mathrm{up}}')$.}
\vspace{0.2cm}\hrule\vspace{0.2cm}
\If{$\|\hat{\boldsymbol{b}}\|_\infty \le \textftol$}{
    \Return $(\mathsf{Ftol}, \mathrm{True}, \mathrm{True}, I_{\mathrm{up}})$ \tcp*{converged on residual}
}
\If{$\neg S_{\mathrm{LS}}$}{
    \If{$\|\hat{\boldsymbol{b}}\|_\infty \le \textstag$}{
        \Return $(\mathsf{NoiseFloor}, \mathrm{True}, \mathrm{True}, I_{\mathrm{up}})$\;
    }
    \Return $(\mathsf{LinesearchFailed}, \mathrm{False}, \mathrm{True}, I_{\mathrm{up}})$\;
}
\If{$\hat\Phi > F_{\mathrm{div}} \Phi$}{
    $I_{\mathrm{up}} \gets I_{\mathrm{up}} + 1$;\quad
    \If{$I_{\mathrm{up}} \ge N_{\mathrm{inc}}$}{
        \Return $(\mathsf{MeritDivergence}, \mathrm{False}, \mathrm{True}, I_{\mathrm{up}})$\;
    }
}\Else{ $I_{\mathrm{up}} \gets 0$\; }
\If{$L \le \textxtol$}{
    \If{$\|\hat{\boldsymbol{b}}\|_\infty \le \textstag$}{
        \Return $(\mathsf{NoiseFloor}, \mathrm{True}, \mathrm{True}, I_{\mathrm{up}})$\;
    }
    \Return $(\mathsf{Xtol}, \mathrm{False}, \mathrm{True}, I_{\mathrm{up}})$\;
}
\If{$i = N_{\mathrm{N}}$}{
    \If{$\|\hat{\boldsymbol{b}}\|_\infty \le \textstag$}{
        \Return $(\mathsf{NoiseFloor}, \mathrm{True}, \mathrm{True}, I_{\mathrm{up}})$\;
    }
    \Return $(\mathsf{IterationCap}, \mathrm{False}, \mathrm{True}, I_{\mathrm{up}})$\;
}
\Return $(\mathsf{InProgress}, \mathrm{False}, \mathrm{False}, I_{\mathrm{up}})$\;
\end{algorithm}

\noindent
In the implementation $I_{\mathrm{up}}$ is not returned explicitly but carried
in the \code{inc\_count} field of a mutable \code{SafeIterationState}
struct that the same routine updates in place; the box treats it as a
returned value for mathematical clarity.

\paragraph{Why six exit reasons?}  The safeguard logic distinguishes
states that classical Newton implementations conflate, because in the
context of the present cascade \emph{the choice of reason changes what
happens next}.
\begin{itemize}
\item $\mathsf{Ftol}$ and $\mathsf{InitialFtol}$ are unambiguous successes;
the caller (Algorithm~\ref{alg:B} or~\ref{alg:O}) proceeds to Stage~II or
exports results.
\item $\mathsf{NoiseFloor}$ is a \emph{conditional} success: the residual
sits below $\textstag$, which is itself far above floating-point noise but
well below physical interest. The caller treats this as a success
intentionally (Section~\ref{sec:falseEquiv} discusses the conditions under
which this classification is and is not safe).
\item $\mathsf{LinesearchFailed}$ and $\mathsf{MeritDivergence}$ both trigger
the Picard fallback (Algorithm~\ref{alg:B}). They differ in cause:
$\mathsf{LinesearchFailed}$ says the Newton direction itself is useful in
direction but cannot find a length that descends; $\mathsf{MeritDivergence}$
says the Newton direction is actively pushing the iterate uphill.
\item $\mathsf{Xtol}$ flags an iterate that has stopped moving even though
the residual is non-trivial; the orchestrator treats this as failure unless
the residual has slipped below $\textstag$ (in which case it is silently
relabelled as $\mathsf{NoiseFloor}$).
\item $\mathsf{IterationCap}$ is the budget-exhausted failure. As with
$\mathsf{Xtol}$, a sufficiently small residual converts it to
$\mathsf{NoiseFloor}$.
\item $\mathsf{LinearSolveNaN}$ aborts the loop unconditionally; the caller
falls back to Picard, which may produce a non-singular linearisation.
\end{itemize}
The asymmetry between failures that promote to $\mathsf{NoiseFloor}$ and
failures that do not is intentional: a stalled iterate \emph{with a small
residual} is a converged iterate that simply cannot be improved further by
the inner solver, whereas $\mathsf{LinesearchFailed}$ and
$\mathsf{MeritDivergence}$ never appear with a small residual---they are
signals from the iteration map itself, not from how close the iterate is to
the root.

\paragraph{The role of $F_{\mathrm{div}}$ and $N_{\mathrm{inc}}$ together.}
A single merit-expansion event is not a divergence: Newton near a saddle can
take one or two steps that increase the merit before the quadratic basin
captures the iterate. Requiring $N_{\mathrm{inc}}$ \emph{consecutive}
expansions, each by at least a factor $F_{\mathrm{div}}$, separates this
transient overshooting from a genuine repeated divergence. Choosing
$N_{\mathrm{inc}}$ small ($\le 2$) makes the guard trip almost immediately
on hard problems and forces an early Picard hand-off; choosing it large
($\ge 5$) lets Newton recover from non-trivial overshoots but at the cost of
diverging further before failing.

% ============================================================
\section{The Picard Fallback}
\label{sec:picard}

When the safeguarded Newton loop terminates without converging, the
orchestration does not give up: it transitions to a Picard iteration map
and then re-engages Newton.

\begin{algorithm}[H]
\caption{$\mathrm{PicardFallback}$: the Stage~I cascade.}
\label{alg:B}
\SetAlgoLined
\textbf{Parameters:} $N_{\mathrm{N}}, N_{\mathrm{P}}, \textftol, \textstag, $ all Newton/line-search parameters.\\
\KwIn{Initial guess $\boldsymbol{x}^0 \in \mathcal{X}_h$ and a backup copy $\boldsymbol{x}^0_{\mathrm{backup}}$.}
\KwOut{$(S, \boldsymbol{x}^\star)$ where $S$ is the overall success flag.}
\vspace{0.2cm}\hrule\vspace{0.2cm}
configure $\boldsymbol{\pi}_h \gets 0$ \tcp*{ASGS-mode residual and Jacobian}
$(\boldsymbol{x}^\star, S_{\mathrm{N}}, \dots, r_{\mathrm{N}}) \gets \mathrm{ExactNewtonPipeline}(\boldsymbol{x}^0,\ \mathrm{ExactNewton\ Jacobian})$\;
\If{$S_{\mathrm{N}} = \mathrm{True}$}{
    \Return $(\mathrm{True}, \boldsymbol{x}^\star)$ \tcp*{Newton sufficed; no fallback}
}

restore $\boldsymbol{x}^0 \gets \boldsymbol{x}^0_{\mathrm{backup}}$\;
$(\boldsymbol{x}^\star, S_{\mathrm{P}}, \dots, r_{\mathrm{P}}) \gets \mathrm{ExactNewtonPipeline}(\boldsymbol{x}^0,\ \mathrm{Picard\ Jacobian})$\;
\If{any component of $\boldsymbol{x}^\star$ is non-finite}{
    \Return $(\mathrm{False}, \boldsymbol{x}^0_{\mathrm{backup}})$ \tcp*{catastrophic state}
}
\If{$S_{\mathrm{P}} = \mathrm{True}$}{
    \Return $(\mathrm{True}, \boldsymbol{x}^\star)$\;
}

\tcc{Re-engage Newton from wherever Picard left the iterate.}
$(\boldsymbol{x}^\star, S_{\mathrm{N2}}, \dots, r_{\mathrm{N2}}) \gets \mathrm{ExactNewtonPipeline}(\boldsymbol{x}^\star,\ \mathrm{ExactNewton\ Jacobian})$\;
\Return $(S_{\mathrm{N2}} \vee S_{\mathrm{P}},\ \boldsymbol{x}^\star)$\;
\end{algorithm}

\paragraph{Why Picard works where Newton fails.}  The Newton Jacobian for
the porous Navier--Stokes system contains the linearised convective term
$\boldsymbol{u}^k \cdot \nabla \delta\boldsymbol{u}
+ \delta\boldsymbol{u} \cdot \nabla \boldsymbol{u}^k$ and the linearised
Forchheimer term $\partial\sigma/\partial\boldsymbol{u}$. Both of these
become large in the difficult regimes, and the Jacobian can develop a near-null
direction where a small change in $\boldsymbol{u}$ produces an outsize change
in the residual. Newton near such a cusp produces a step that diverges. The
Picard Jacobian drops the dependence of $\boldsymbol{u}\cdot\nabla\boldsymbol{u}$
and of $\sigma$ on the increment $\delta\boldsymbol{u}$, retaining only the
``frozen'' coefficient form. The resulting iteration map is a contraction in
a wider region than Newton's quadratic basin, at the cost of linear (rather
than quadratic) convergence rate. Picard does not converge to within
$\textftol$ in general---but it does push the iterate into a region where
the subsequent Newton retry has a fighting chance.

\paragraph{Triggers.}  The original algorithm box in earlier drafts of this
document restricted the Picard entry to two exit reasons. The current
implementation is more permissive: \emph{any} non-converged Newton exit
triggers the fallback, including $\mathsf{Xtol}$ stagnation,
$\mathsf{IterationCap}$, the linear-solver exception path, and any uncaught
Julia exception during the Newton call. This is appropriate because all of
these states leave $\boldsymbol{x}^0$ outside the basin, and a Picard pass
on top costs the same regardless of which symptom triggered the entry.

\paragraph{The catastrophic-state guard.}  Picard, while contractive in
benign regimes, is not bounded in general. After the Picard call, the
orchestration checks $\mathrm{any}(!\mathrm{isfinite}, \boldsymbol{x}^\star)$.
A NaN or Inf in the iterate is a hard signal that even Picard could not
keep the model evaluable, and the fallback restores the original backup and
returns failure. Without this guard the subsequent Newton call would crash
on its first residual evaluation.

\paragraph{Connection to Algorithm~\ref{alg:O}.}  Stage~I of the
orchestration is exactly one call to $\mathrm{PicardFallback}$. The
``Picard-smoothed'' state produced by line~8 of Algorithm~\ref{alg:B} is
$U_h^{\mathrm{ASGS}}$ in Algorithm~\ref{alg:O}, and it is the seed for the
OSGS staggered loop.

% ============================================================
\section{OSGS Staggered Relaxation}
\label{sec:osgs}

The orthogonal sub-grid scale formulation augments the discrete weak form
with a stabilisation term that depends on the $L^2$ projection
$\boldsymbol{\pi}_h$ of the strong residual. To avoid carrying the
Fr\'echet derivative of $\Pi_h$ in the global Jacobian, the algorithm
treats the pair $(U_h, \boldsymbol{\pi}_h)$ by block-coordinate
relaxation: freeze $\boldsymbol{\pi}_h$, solve the inner nonlinear primary
system for $U_h$, then re-project to get the next $\boldsymbol{\pi}_h$.

\subsection{Algorithm C}

\begin{algorithm}[H]
\caption{$\mathrm{OSGSFractionalRelaxation}$: staggered fixed point on $(U_h, \boldsymbol{\pi}_h)$.}
\label{alg:C}
\SetAlgoLined
\textbf{Parameters:} $N_{\mathrm{OSGS}}, N_{\mathrm{inner}}, m_{\mathrm{w}}, \tau_{\mathrm{w}}, \textftol, \tau_{\mathrm{OSGS}}, \tau_{\mathrm{proj}}, \textstag, \mathsf{Mode}_{\mathrm{drift}}, \mathsf{Mode}_{\mathrm{stop}}$.\\
\KwIn{Stage--I converged state $U_h^{\mathrm{ASGS}}$; residual operators $(\mathcal{R}_u, \mathcal{R}_p)$; pre-factored mass solvers $(M_u^{-1}, M_p^{-1})$; optional Anderson accelerators.}
\KwOut{$(S, \boldsymbol{\pi}_{h,u}, \pi_{h,p}, \dots)$.}
\vspace{0.2cm}\hrule\vspace{0.2cm}
$U_h^{0} \gets U_h^{\mathrm{ASGS}}$\;
$\boldsymbol{\pi}_{h,u}^{0} \gets \Pi_{V_{\mathrm{free}}}(\mathcal{R}_u(U_h^{0}))$;\quad
$\pi_{h,p}^{0} \gets \Pi_{Q_{\mathrm{free}}}(\mathcal{R}_p(U_h^{0}))$\;
\For{$m = 1, \dots, N_{\mathrm{OSGS}}$}{
    set effective inner tolerance $\ftolinner{m} \gets (m \le m_{\mathrm{w}}) \,?\, \max(\textftol, \tau_{\mathrm{w}}) : \textftol$\;
    fix $\boldsymbol{\pi}_h \gets \boldsymbol{\pi}_h^{m-1}$;\quad
    $(U_h^{m}, S_m^{\mathrm{inner}}, \dots, r_m^{\mathrm{inner}}) \gets \mathrm{ExactNewtonPipeline}(U_h^{m-1})$ with cap $N_{\mathrm{inner}}$ and tolerance $\ftolinner{m}$\;
    \If{$r_m^{\mathrm{inner}} \in \{\mathsf{LinesearchFailed}, \mathsf{MeritDivergence}, \mathsf{LinearSolveNaN}\}$}{
        \Return $(\mathrm{False}, \boldsymbol{\pi}_{h,u}^{m-1}, \pi_{h,p}^{m-1}, \dots)$ \tcp*{abort: structural OSGS divergence}
    }

    \tcc{State drift, in the configured metric.}
    \If{$\mathsf{Mode}_{\mathrm{drift}} = \code{L2\_mass}$}{
        $d_{\mathrm{state}}^m \gets \max\!\bigl(\|\boldsymbol{u}_h^m - \boldsymbol{u}_h^{m-1}\|_{L^2(\Omega)},\ \|p_h^m - p_h^{m-1}\|_{L^2(\Omega)}\bigr)$\;
    }\Else{
        $d_{\mathrm{state}}^m \gets \|\boldsymbol{X}_h^m - \boldsymbol{X}_h^{m-1}\|_{\ell^\infty}$ over the global DOF vector\;
    }

    re-project: $\boldsymbol{\pi}_{h,u}^{m} \gets \Pi_{V_{\mathrm{free}}}(\mathcal{R}_u(U_h^m))$;\quad
    $\pi_{h,p}^{m} \gets \Pi_{Q_{\mathrm{free}}}(\mathcal{R}_p(U_h^m))$\;

    \tcc{Optional block-wise Anderson acceleration on each projection.}
    \If{accelerator present}{
        $\boldsymbol{\pi}_{h,u}^{m} \gets \mathrm{Anderson}_u(\boldsymbol{\pi}_{h,u}^{m-1}, \boldsymbol{\pi}_{h,u}^{m})$;\quad
        $\pi_{h,p}^{m} \gets \mathrm{Anderson}_p(\pi_{h,p}^{m-1}, \pi_{h,p}^{m})$\;
    }

    $d_\pi^m \gets \max\!\bigl(
        \sqrt{(\Delta\boldsymbol{\pi}_{h,u}^m)^\top M_u (\Delta\boldsymbol{\pi}_{h,u}^m)},\
        \sqrt{(\Delta\pi_{h,p}^m)^\top M_p (\Delta\pi_{h,p}^m)} \bigr)$\;

    \tcc{Convergence predicate, noise-floor protected.}
    $\tau_{\mathrm{OSGS}}^{\mathrm{eff}} \gets \max(\tau_{\mathrm{OSGS}}, \textstag)$;\quad
    $\tau_{\mathrm{proj}}^{\mathrm{eff}} \gets \max(\tau_{\mathrm{proj}}, \textstag)$\;
    $S_{\mathrm{state}} \gets (d_{\mathrm{state}}^m \le \tau_{\mathrm{OSGS}}^{\mathrm{eff}})$;\quad
    $S_{\mathrm{proj}}  \gets (d_\pi^m       \le \tau_{\mathrm{proj}}^{\mathrm{eff}})$\;
    $S_{\mathrm{conv}} \gets \begin{cases}
        S_{\mathrm{state}} & \mathsf{Mode}_{\mathrm{stop}} = \code{state\_drift}, \\
        S_{\mathrm{proj}}  & \mathsf{Mode}_{\mathrm{stop}} = \code{projection\_drift}, \\
        S_{\mathrm{state}} \wedge S_{\mathrm{proj}} & \mathsf{Mode}_{\mathrm{stop}} = \code{both}.
    \end{cases}$\;
    \If{$m \ge 2$ \textnormal{\textbf{and}} $\|\boldsymbol{X}_h^m - \boldsymbol{X}_h^{m-1}\|_{\ell^\infty} \le \textstag$}{
        $S_{\mathrm{conv}} \gets \mathrm{True}$ \tcp*{absolute noise-floor short circuit}
    }
    \If{$S_{\mathrm{conv}}$}{
        \textbf{[}\,if MMS enabled, run $\mathrm{VerifyMMSPlateau}$ (Alg.~\ref{alg:D})\,\textbf{]}\;
        \Return $(\mathrm{True}, \boldsymbol{\pi}_{h,u}^m, \pi_{h,p}^m, \dots)$\;
    }
}
\Return $(\mathrm{False}, \boldsymbol{\pi}_{h,u}^{N_{\mathrm{OSGS}}}, \pi_{h,p}^{N_{\mathrm{OSGS}}}, \dots)$\;
\end{algorithm}

\subsection{The bootstrap and why the ASGS state matters}

The first sub-grid projection $\boldsymbol{\pi}_h^0$ is the $L^2$ projection
of the strong residual evaluated at $U_h^{\mathrm{ASGS}}$, not at the zero
initial guess. There are two reasons for this.

First, when the ASGS system has converged, $\mathcal{R}_u(U_h^{\mathrm{ASGS}})$
and $\mathcal{R}_p(U_h^{\mathrm{ASGS}})$ are small in $L^2$, so
$\boldsymbol{\pi}_h^0$ is already close to its fixed-point value.
The staggered loop is then a perturbation around an existing solution
rather than a global search.

Second, an ASGS-converged primary field gives the inner Newton solve
at $m = 1$ an initial residual that is small in absolute terms, so the
Newton iteration usually converges within a handful of inner iterations.
Starting OSGS from $\boldsymbol{x}^0 = 0$, by contrast, would force the
solver to chase both the primary nonlinearity and the projection coupling
simultaneously, which empirically diverges in stiff regimes.

\subsection{Why projection happens on the unconstrained spaces}

The projection target $V_{\mathrm{free}}, Q_{\mathrm{free}}$ has no
Dirichlet constraints. If the projection were performed onto the
constrained trial spaces $U, P$, the boundary residual---which is
non-zero in general---would be assigned to the projection's boundary DOFs,
producing an artefact of order $O(1)$ along $\partial\Omega$. That artefact
then enters the stabilisation term, where it pollutes the convergence rate of
the FE method to $O(1)$ rather than the expected $O(h^{k+1})$, destroying any
MMS rate verification. See \code{theory/paper-code-divergences.md}.

\subsection{Mass-matrix pre-factoring}

Each projection $\Pi_{V_{\mathrm{free}}}, \Pi_{Q_{\mathrm{free}}}$ is the
solution of an SPD linear system with the mass matrices $M_u, M_p$. Because
the spaces $V_{\mathrm{free}}, Q_{\mathrm{free}}$ do not change between
outer iterations, the matrices and their LU factorisations are assembled
\emph{once} (Algorithm~\ref{alg:O}, line~16). At every outer iteration the
projection cost is then one forward and one back substitution per block;
the cubic factorisation cost is amortised over the full $N_{\mathrm{OSGS}}$
budget.

\subsection{Warmup phase}

For the first $m_{\mathrm{w}}$ outer iterations the inner Newton tolerance is
relaxed to $\max(\textftol, \tau_{\mathrm{w}})$. The motivation is purely
economic: while $\boldsymbol{\pi}_h^0$ is still being computed off a
not-yet-stationary residual, driving the inner solve to a tight $\textftol$
``polishes'' the iterate against a projection that will be replaced
immediately. Relaxing the inner tolerance during this phase trades a tiny
amount of accuracy in the early iterates for a noticeable cut in inner-Newton
budget. Once $m > m_{\mathrm{w}}$ the inner tolerance reverts to the global
$\textftol$ and the staggered loop tightens.

A practical default has $m_{\mathrm{w}} \in \{1,2,3\}$ and $\tau_{\mathrm{w}}
= 10\cdot\textftol$ or so. Increasing $m_{\mathrm{w}}$ further wastes
fewer Newton iterations early but risks the staggered fixed point
``locking'' on a coarse projection.

\subsection{State drift: \texorpdfstring{$\ell^\infty$ vs continuous $L^2$}{Linf vs L2}}
\label{sec:drift-modes}

The state drift between consecutive outer iterations
\begin{equation}
    d_{\mathrm{state}}^m
    = \begin{cases}
        \max\!\bigl(
            \|\boldsymbol{u}_h^m - \boldsymbol{u}_h^{m-1}\|_{L^2(\Omega)},\
            \|p_h^m - p_h^{m-1}\|_{L^2(\Omega)} \bigr)
            & \mathsf{Mode}_{\mathrm{drift}} = \code{L2\_mass}, \\[1mm]
        \|\boldsymbol{X}_h^m - \boldsymbol{X}_h^{m-1}\|_{\ell^\infty}
            & \mathsf{Mode}_{\mathrm{drift}} = \code{Linf},
    \end{cases}
\end{equation}
admits two metrics with subtly different meanings.

\paragraph{The \texttt{Linf} mode.}  This is the cheapest possible drift:
the raw vector-space $\ell^\infty$ norm of the concatenated free-DOF array.
It costs one BLAS sweep and no quadrature. The pathology, however, is that
the same array stores velocity and pressure DOFs interleaved by the
multi-field layout, so a velocity DOF expressed in (dimensional) units of
$[\mathrm{m/s}]$ sits beside a pressure DOF in $[\mathrm{Pa}]$. The
$\ell^\infty$ norm is then dominated by whichever field happens to have the
larger numerical magnitude---a scale-mixing artefact, not a physical drift
measurement. In regimes where pressure and velocity are well balanced this
mode works adequately and is fast. In regimes where one field dominates,
the drift signal can become essentially blind to the other.

\paragraph{The \texttt{L2\_mass} mode.}  Computes two independent continuous
$L^2(\Omega)$ norms via quadrature, one for each block, and takes their
maximum. This is the recommended default (and the invariant codified in
\code{CLAUDE.md}): velocity and pressure are measured in their own
natural functional norms, with no possibility of cross-pollution. The
cost is two mass-weighted quadratures per outer iteration, dominated in
practice by the inner Newton solve itself.

\subsection{Projection drift, in the mass-weighted inner product}

The projection drift
\begin{equation}
    d_\pi^m
    = \max\!\Bigl(
        \sqrt{(\Delta\boldsymbol{\pi}_{h,u}^m)^\top M_u (\Delta\boldsymbol{\pi}_{h,u}^m)},\
        \sqrt{(\Delta\pi_{h,p}^m)^\top M_p (\Delta\pi_{h,p}^m)} \Bigr)
\end{equation}
is the natural distance in the inner product that defines the $L^2$
projection itself. Because the projection is the minimiser of a quadratic
in $M$, the appropriate measure of progress toward the fixed point is the
$M$-weighted norm; the Euclidean DOF norm would conflate it with
inverse-mass scaling and report misleading drift values on graded meshes.
The factorisations $M_u^{-1}, M_p^{-1}$ are not needed for this norm; only
the matrix-vector products $M_u \Delta\pi, M_p \Delta\pi$ are required, and
those reuse the same assembled sparse matrices used in the projection step.

\subsection{Choosing the stopping mode}

The user selects $\mathsf{Mode}_{\mathrm{stop}}$ between three options:
\begin{itemize}
\item \code{state\_drift}: terminate when $d_{\mathrm{state}}^m
\le \tau_{\mathrm{OSGS}}^{\mathrm{eff}}$. Use when the primary fields are
the only quantities of interest and the projection is a means to an end.
\item \code{projection\_drift}: terminate when $d_\pi^m \le
\tau_{\mathrm{proj}}^{\mathrm{eff}}$. Use when sub-grid scale accuracy is
the goal and one expects the primary field to be approximately stationary
already.
\item \code{both} (default): require both predicates. The most defensive
option and the one that aligns with the mathematical definition of the
fixed point.
\end{itemize}
In all three cases an additional absolute short-circuit fires when the raw
$\ell^\infty$ DOF drift drops below $\textstag$ for $m \ge 2$; this protects
against limit-cycle oscillations in which the iterate ping-pongs between
two nearby states without either drift ever cleanly dropping below
$\tau_{\mathrm{OSGS}}$.

\subsection{Block-wise Anderson acceleration}
\label{sec:anderson}

For high-$\mathrm{Re}$, high-$\mathrm{Da}$ flows the staggered fixed point
on $\boldsymbol{\pi}_h$ can stall in slow limit-cycle oscillations. Anderson
acceleration extrapolates a new $\boldsymbol{\pi}_h^m$ from the history of
the last $m$ fixed-point residuals by solving a small mass-weighted
least-squares problem
\begin{equation}
    \min_\gamma\ \|(\Delta F)\gamma - f_k\|_M^2,
\end{equation}
where $\Delta F$ is the matrix of residual increments and the inner product
$\|\cdot\|_M^2 = (\cdot)^\top M (\cdot)$ uses the appropriate mass matrix.
Two architectural choices make this robust in the present setting:

\paragraph{Block isolation.}  The accelerator operates independently on
$\boldsymbol{\pi}_{h,u}$ and $\pi_{h,p}$, each with its own mass matrix
($M_u$ and $M_p$ respectively). Combining the blocks into a single
extrapolation problem would force a cross-physical mixing that the OSGS
formulation goes out of its way to avoid in the first place.

\paragraph{Mass-weighted normal equations.}  The least-squares system
uses the same $L^2$ inner product that defines the underlying projection,
not the Euclidean dot product of DOF vectors. This is the natural way to
make the extrapolation respect mesh grading.

\paragraph{Safety fallback.}  The implementation contains a sanity check:
if the extrapolated state $x_{\mathrm{next}}$ lies further than $10$ times
the unextrapolated residual from the raw fixed-point candidate $g_k$, the
extrapolation is discarded and the iteration falls back to a relaxed
Picard step $x_k + \beta f_k$. The factor $10$ is a conservative bound
chosen so that the safety net does not interfere with legitimate
Anderson acceleration on well-behaved problems while still catching
catastrophic over-extrapolation in degenerate cases.

\subsection{The reaction-projection trim}

When the configured \code{experimental\_reaction\_mode} is \code{"standard"}
\emph{and} the reaction law is \code{ConstantSigmaLaw}, the projected
residual omits the $\sigma\boldsymbol{u}$ term entirely
(\code{src/stabilization/projection.jl}:26--40). Mathematically this is
the observation that for constant $\sigma$ the $L^2$ projection of
$\sigma\boldsymbol{u}$ onto the FE space is itself, so the contribution
to $\Pi_h(\mathcal{R}_u)$ exactly cancels in the OSGS residual. Skipping it
saves an assembly per outer iteration without altering the converged state.
For non-constant reaction laws (e.g.\ Forchheimer--Ergun) the trim is
invalid and the projection-policy sanitiser rejects this option.

\subsection{Failure of the inner Newton}

The inner Newton call inside Algorithm~\ref{alg:C} runs on the same
$\mathrm{ExactNewtonPipeline}$ used in Algorithm~\ref{alg:A}, but with the
iteration cap reduced to $N_{\mathrm{inner}}$ (typically a small number,
e.g.\ $3$--$5$) and the tolerance possibly relaxed during warmup. When the
inner exit reason is one of $\mathsf{LinesearchFailed}$,
$\mathsf{MeritDivergence}$, or $\mathsf{LinearSolveNaN}$, the staggered
outer loop aborts immediately: these reasons reflect structural failure of
the inner iteration map, and continuing the outer relaxation would only
re-feed the broken state into another projection. Other exit reasons
(notably $\mathsf{IterationCap}$ from exhausting $N_{\mathrm{inner}}$
without converging to the effective inner tolerance) are treated as
progress: the next outer iteration simply re-engages with the updated
projection.

% ============================================================
\section{MMS Plateau Verification}
\label{sec:mms}

\subsection{Numerical convergence versus discretisation convergence}

The inner Newton stops when $\|\boldsymbol{b}\|_\infty \le \textftol$.
That is convergence in the \emph{algebraic} sense: the iterate is close to
the root of the discrete system. It says nothing about whether the discrete
system itself is close to the true PDE solution. The Method of Manufactured
Solutions (MMS) provides exactly that second test: it injects an exact
analytical $(\boldsymbol{u}_{\mathrm{ex}}, p_{\mathrm{ex}})$, recomputes
the appropriate source terms, and measures
$\|U_h - U_{\mathrm{ex}}\|$ in physical norms. The expected discretisation
behaviour is $O(h^{k+1})$ in the $L^2$ norms (and one order weaker in the
$H^1$-semi norm).

For a single mesh, plateau verification asks a different question: once
the iterate is close enough to the root that further iterations do not
move the FE error, the FE error has \emph{plateaued} at its discretisation
floor. Detecting this plateau certifies that the algebraic convergence
criterion was tight enough to expose the discretisation error.

\subsection{Algorithm D}

\begin{algorithm}[H]
\caption{$\mathrm{VerifyMMSPlateau}$: relative-change plateau test.}
\label{alg:D}
\SetAlgoLined
\textbf{Parameters:} $N_{\mathrm{plat}}, \tau_{\mathrm{err}}, \varepsilon_{u,L^2}, \varepsilon_{u,H^1}, \varepsilon_{p,L^2}$.\\
\KwIn{Current iterate $U_h^m$, exact reference $U_{\mathrm{ex}}$, history $\mathcal{H}$ of past error tuples.}
\KwOut{$(S_{\mathrm{plat}})$.}
\vspace{0.2cm}\hrule\vspace{0.2cm}
compute $E_{u,L^2}^m = \|\boldsymbol{u}_h^m - \boldsymbol{u}_{\mathrm{ex}}\|_{L^2}$,
        $E_{u,H^1}^m = |\boldsymbol{u}_h^m - \boldsymbol{u}_{\mathrm{ex}}|_{H^1}$,
        $E_{p,L^2}^m = \|p_h^m - p_{\mathrm{ex}}\|_{L^2}$\;
append to $\mathcal{H}$; let $E^{m-1}_\bullet$ be the previous entry\;
$r_{u,L^2} \gets |E_{u,L^2}^m - E_{u,L^2}^{m-1}| / \max(E_{u,L^2}^m,E_{u,L^2}^{m-1}, \varepsilon_{u,L^2})$\;
$r_{u,H^1} \gets |E_{u,H^1}^m - E_{u,H^1}^{m-1}| / \max(E_{u,H^1}^m,E_{u,H^1}^{m-1}, \varepsilon_{u,H^1})$\;
$r_{p,L^2} \gets |E_{p,L^2}^m - E_{p,L^2}^{m-1}| / \max(E_{p,L^2}^m,E_{p,L^2}^{m-1}, \varepsilon_{p,L^2})$\;
$r_{\max} \gets \max(r_{u,L^2},r_{u,H^1},r_{p,L^2})$\;
\If{$r_{\max} \le \tau_{\mathrm{err}}$}{
    $I_{\mathrm{pass}} \gets I_{\mathrm{pass}} + 1$\;
}\Else{
    $I_{\mathrm{pass}} \gets 0$\;
}
\If{$I_{\mathrm{pass}} \ge N_{\mathrm{plat}}$}{
    \Return $\mathrm{True}$ \tcp*{plateau satisfied}
}
\Return $\mathrm{False}$ \tcp*{continue extending}
\end{algorithm}

\subsection{Three norms, three floors}

The verifier tracks \emph{three} error quantities, not two: $L^2$ of
velocity, $L^2$ of pressure, and the $H^1$-semi-norm
$|\boldsymbol{u}_h - \boldsymbol{u}_{\mathrm{ex}}|_{H^1}
= \sqrt{\int_\Omega \nabla(\boldsymbol{u}_h - \boldsymbol{u}_{\mathrm{ex}}) :
\nabla(\boldsymbol{u}_h - \boldsymbol{u}_{\mathrm{ex}})\,\mathrm{d}\Omega}$.
The third norm matters because pressure $L^2$ and velocity $L^2$ can
plateau while the velocity gradient---which is what the convective and
viscous terms actually feel---still drifts. Without the $H^1$-semi norm a
silent gradient drift would be invisible to the plateau test.

Each ratio is taken with respect to its own per-norm floor
$\varepsilon_\bullet$. Without these floors, when an error first drops to
$10^{-15}$ the ratio between consecutive iterations becomes numerically
meaningless. The floors are configured high enough above floating-point
noise to remain stable on fine meshes, but low enough that they do not
artificially declare convergence for genuinely coarse runs.

\subsection{Consecutive passes and budget exhaustion}

The plateau is declared only after $N_{\mathrm{plat}}$ \emph{consecutive}
passes, and the counter resets to zero on the first failure. This catches
intermittent drift---e.g.\ a sequence
$(r_{\max}^k)_k = (10^{-4}, 10^{-3}, 10^{-4}, \ldots)$
that looks safe sample-by-sample but actually has not stabilised.
A typical choice is $N_{\mathrm{plat}} \in \{2,3\}$.

\subsection{Two entry points and how they differ}

The verifier is invoked from two distinct contexts.

\paragraph{ASGS path.}  After Stage~I converges and the method is ASGS, the
orchestration enters a dedicated extension loop of length $K_{\mathrm{extra}}$
(\code{solve\_system}, lines~258--337). Each iteration calls
$\mathrm{ExactNewtonPipeline}$ with $N_{\mathrm{N}} = 1$ followed by
$\mathrm{VerifyMMSPlateau}$. The loop exits when the plateau passes or the
budget is exhausted; either way the iterate is returned. The ``one inner
Newton step then verify'' cadence is deliberate: it gives the plateau test
a chance to observe how a single extra Newton step changes the FE error.

\paragraph{OSGS path.}  Within $\mathrm{OSGSFractionalRelaxation}$
(Algorithm~\ref{alg:C}), the verifier runs only after the staggered
predicate $S_{\mathrm{conv}}$ has fired, not at every outer iteration.
The first successful staggered iteration triggers a fresh MMS error history;
subsequent successful iterations append to it and test for the plateau. The
budget inflation $K_{\mathrm{extra}}$ allows the staggered loop to keep
producing additional converged iterates after first hitting the fixed point.
This entry point is gated rather than per-iteration because the FE error is
defined relative to the discrete system, and that system itself is moving
between staggered iterations until $\boldsymbol{\pi}_h$ stabilises.

% ============================================================
\section{The MMS Sweep Harness}
\label{sec:harness}

The discussion above applies to a single \code{run\_simulation} call.
The MMS sweep harness under \code{test/}\allowbreak\code{extended/}\allowbreak\code{ManufacturedSolutions/}
wraps many \code{run\_simulation} runs in a triple loop over Reynolds number
$\mathrm{Re}$,
Damk\"ohler number $\mathrm{Da}$, and mesh resolution $h$, and reports
convergence-rate tables across the resulting sweep. Because the harness sees
many cases of widely varying difficulty, it injects three per-case
adaptations that the main \code{run\_simulation} flow deliberately does not
attempt. We describe them here in one place so that the reader does not
mistake them for parts of the orchestration.

\subsection{Adaptive \texorpdfstring{$\textftol$}{ftol}}
\label{sec:harness-ftol}

Two error scales matter in an FE method. The \emph{algebraic error} is the
distance between the iterate and the discrete root, governed by $\textftol$.
The \emph{discretisation error} is the distance between the discrete root
and the true PDE solution, which behaves as $O(h^{k_v+1})$ for the chosen
velocity element order. Setting $\textftol$ much tighter than $O(h^{k_v+1})$
wastes Newton iterations driving the algebraic residual far below the
discretisation floor---adjustments that cannot improve the discretisation
error. Setting $\textftol$ much looser than $O(h^{k_v+1})$ lets the
algebraic error dominate and pollutes the observed convergence rate.

The harness therefore re-derives $\textftol$ per case:
\begin{equation}
    \label{eq:dynamic-ftol}
    \textftol^{\,\mathrm{dyn}}
    = \max\!\Bigl(\textftol^{\,\mathrm{static}},\
        \min(C_{\mathrm{ceil}},\ C_{\mathrm{sf}} \cdot h^{k_v+1}) \Bigr).
\end{equation}
The right-most term ties the tolerance to the discretisation scale via a
spatial safety factor $C_{\mathrm{sf}}$ (a typical choice is $10^{-2}$,
i.e.\ the algebraic residual must drop one to two orders of magnitude
below the predicted discretisation error). The ceiling $C_{\mathrm{ceil}}$
prevents trivially loose tolerances on absurdly coarse meshes. The outer
$\max$ prevents the formula from descending below the user-supplied
$\textftol^{\,\mathrm{static}}$, which serves as a sanity bound. The
adaptive $\textftol^{\,\mathrm{dyn}}$ is then plugged into every
\code{SafeNewtonSolver} the harness constructs for that case
(\code{run\_test.jl}:440--446, 470--471).

The main \code{run\_simulation} driver does not use this scaling because
in production work the user typically picks one mesh and one tolerance
deliberately, and an adaptive rule would silently override that choice.

\subsection{Adaptive noise floor}
\label{sec:harness-noise}

Linear-system conditioning grows as $\mathcal{O}(h^{-2})$ for elliptic
operators and is further amplified by convective and reactive terms. A
static $\textstag$ that is appropriate on a coarse mesh will be far above
the genuine round-off floor on a fine mesh, causing the
$\mathsf{NoiseFloor}$ classification of Algorithm~\ref{alg:A3} to fire
prematurely. The harness consequently scales
\begin{equation}
    \textstag^{\,\mathrm{dyn}}
    = \min\!\Bigl(\textstag,\ \max(n^{\min},\ n^{\mathrm{base}} \cdot N^2 \cdot \max(1, \mathrm{Re}))\Bigr),
\end{equation}
clipped above by the user-supplied $\textstag$ and below by an absolute
floor $n^{\min}$. The factor $N^2$ tracks the elliptic conditioning growth
(with $N$ the partition count per direction, so $h \propto 1/N$); the
$\max(1,\mathrm{Re})$ factor accounts for convective amplification.
A final $\textstag^{\,\mathrm{dyn}} \gets
\max(\textstag^{\,\mathrm{dyn}}, n^{\mathrm{sf}}\,\textftol^{\,\mathrm{dyn}})$
clamp preserves the strict ordering $\textstag^{\,\mathrm{dyn}} >
\textftol^{\,\mathrm{dyn}}$, without which Algorithm~\ref{alg:A3} would
collapse the two predicates.

\subsection{Per-case Picard widening}
\label{sec:harness-picard}

In a $(\mathrm{Re}, \mathrm{Da}, h)$ sweep, the easy cases require few
Picard iterations and the hard ones need many. Setting a uniform
$N_{\mathrm{P}}$ across the sweep forces an unhappy choice: large enough
for the hard cases (wastefully slow on the easy ones) or small enough for
the easy ones (insufficient on the hard ones). The harness therefore
\emph{widens} the Picard budget when a stiffness threshold is exceeded
(\code{run\_test.jl}:459--468):
\begin{equation}
    \begin{aligned}
    N_{\mathrm{P}}^{\,\mathrm{local}} &\gets N_{\mathrm{P}}, \\
    \text{if } \mathrm{Re} \ge \mathrm{Re}^{\mathrm{th}}: \quad
    N_{\mathrm{P}}^{\,\mathrm{local}} &\gets \max(N_{\mathrm{P}}^{\,\mathrm{local}}, N_{\mathrm{P}}^{\mathrm{Re}}), \\
    \text{if } \mathrm{Da} \ge \mathrm{Da}^{\mathrm{th}}: \quad
    N_{\mathrm{P}}^{\,\mathrm{local}} &\gets \max(N_{\mathrm{P}}^{\,\mathrm{local}}, N_{\mathrm{P}}^{\mathrm{Da}}).
    \end{aligned}
\end{equation}
$N_{\mathrm{P}}^{\,\mathrm{local}}$ is then plugged into the local
\code{SafeNewtonSolver} for that case. The defaults
$\mathrm{Re}^{\mathrm{th}} = \mathrm{Da}^{\mathrm{th}} = 10^4$ are crude
break-points between regimes where the convective or reactive coupling
becomes the dominant difficulty.

\subsection{Closing remark}

None of the adaptations of Sections~\ref{sec:harness-ftol}--\ref{sec:harness-picard}
fire inside a single \code{run\_simulation} call. The configuration schema
exposes them so that the harness can read them from a shared JSON file,
but the production driver simply uses the static $\textftol$, $\textstag$,
and $N_{\mathrm{P}}$ written in the user's configuration.

% ============================================================
\section{Convergence-Analysis Notes}
\label{sec:falseEquiv}

\subsection{The false-equivalency failsafe}

The single most important invariant of Algorithm~\ref{alg:A3} is the
asymmetry between two exit reasons that, at first glance, look
interchangeable:
\begin{itemize}
\item $\mathsf{IterationCap}$ with $\|\boldsymbol{b}\|_\infty \gg \textstag$:
the iteration ran out of budget while still far from any noise floor.
\item $\mathsf{NoiseFloor}$ with $\|\boldsymbol{b}\|_\infty \le \textstag$:
the iteration may have stopped for any reason, but the residual is below
the noise floor we set for ``effectively converged''.
\end{itemize}
The orchestration treats only the second as success. Coercing the first
into a converged state---``it ran for many iterations and stopped; close
enough''---is forbidden, because doing so would feed an unresolved iterate
into Stage~II of the cascade. In OSGS in particular, an unresolved
$U_h$ produces a stale $\boldsymbol{\pi}_h$ whose subsequent fixed-point
iterations cannot recover.

To make this concrete: imagine a coarse-mesh MMS run in which Newton
stops at $\mathsf{IterationCap}$ with $\|\boldsymbol{b}\|_\infty = 10^{-3}$,
above $\textstag = 10^{-5}$ but well above $\textftol = 10^{-10}$. The
discretisation error on that mesh might be $10^{-4}$. If $\mathsf{IterationCap}$
were silently relabelled as converged, the verifier in
Algorithm~\ref{alg:D} would compute relative-change ratios that reflect a
non-stationary iterate, the plateau test would either oscillate or fail
spuriously, and the reported $O(h^{k+1})$ rate would be corrupted. The
asymmetry in Algorithm~\ref{alg:A3} prevents exactly this.

\subsection{Why $\textstag$ must scale with conditioning}

The same invariant explains why the harness adapts $\textstag$ to the
mesh and to $\mathrm{Re}$ (Section~\ref{sec:harness-noise}). On a fine
mesh the legitimate machine-noise floor lies well below the coarse-mesh
value, and a fixed $\textstag$ would silently rule out states that, while
not at $\textftol$, are genuinely as close to it as the floating-point
arithmetic allows. The scaling
$\textstag \propto N^2 \cdot \max(1, \mathrm{Re})$ tracks the
condition number well enough to keep the $\mathsf{NoiseFloor}$
classification meaningful across the sweep.

\subsection{Why the OSGS stopping mode matters}

Algorithm~\ref{alg:C} accepts three convergence predicates. The default
\code{both} is the strictest and corresponds to the mathematical
definition of the staggered fixed point: $(U_h, \boldsymbol{\pi}_h)$ is
the converged pair when neither moves between outer iterations. The
relaxed modes \code{state\_drift} and \code{projection\_drift} sacrifice
one of the two conditions in exchange for fewer iterations. They are
appropriate only when the user knows that the slacker condition is
already satisfied---e.g.\ because the projection naturally lies in a
regime where it has stabilised quickly, but the primary fields are still
adjusting. Using a single-criterion mode without that knowledge introduces
exactly the false-convergence trap that the noise-floor invariant
forbids: the staggered loop declares victory while one of the two
quantities is still drifting.

% ============================================================
\appendix
\section{Code--Notation Correspondence}
\label{sec:appendix-codes}

Table~\ref{tab:reasons} maps every exit-reason label used by the
algorithm boxes to the string literal returned by the implementation.

\begin{table}[h]
\centering
\caption{Exit reasons in Algorithm~\ref{alg:A3} and their code-side strings.}
\label{tab:reasons}
\small
\begin{tabularx}{\textwidth}{@{} l l X @{}}
\toprule
\textbf{Box label} & \textbf{Code string} & \textbf{Meaning} \\
\midrule
$\mathsf{Ftol}$               & \code{"ftol\_reached"}                 & $\|\boldsymbol{b}\|_\infty \le \textftol$ after a successful line search. \\
$\mathsf{InitialFtol}$        & \code{"initial\_ftol"}                  & Initial guess already satisfies $\textftol$; no Newton step taken. \\
$\mathsf{NoiseFloor}$         & \code{"stagnation\_noise\_floor\_reached"} & $\|\boldsymbol{b}\|_\infty \le \textstag$ on a non-converged exit; treated as success. \\
$\mathsf{LinesearchFailed}$   & \code{"linesearch\_failed"}             & Line-search step shrank to $\alpha \le \alpha_{\min}$ without descent. \\
$\mathsf{MeritDivergence}$    & \code{"merit\_divergence\_escaped"}     & Merit function expanded by $\ge F_{\mathrm{div}}$ for $N_{\mathrm{inc}}$ consecutive steps. \\
$\mathsf{Xtol}$               & \code{"xtol\_stagnation"}               & Step norm $L \le \textxtol$ with residual still above $\textstag$. \\
$\mathsf{IterationCap}$       & \code{"max\_iters\_stagnation"}         & Iteration budget exhausted with residual still above $\textstag$. \\
$\mathsf{LinearSolveNaN}$     & \code{"linear\_solve\_nan"}             & Linear-solve exception or NaN in $\delta\boldsymbol{x}$. \\
$\mathsf{InProgress}$         & \code{""}                                & Sentinel for ``no break this iteration''; never visible to the caller. \\
\bottomrule
\end{tabularx}
\end{table}

The \code{SafeNewtonSolver} struct in \code{src/solvers/nonlinear.jl}
declares all Newton-side numerical parameters explicitly via its
constructor signature; the \code{StabilizationConfig} struct in
\code{src/config.jl} declares the OSGS-side parameters; the
\code{AcceleratorConfig} sub-struct declares the Anderson parameters.
Every parameter listed in Table~\ref{tab:params} corresponds to exactly one
field across these three structures, validated by the \code{validate!}
routine of \code{src/config.jl}; the routine rejects unknown JSON keys
and enforces strict ranges (e.g.\ $0 < c_1 < 1$, $F_{\mathrm{div}} \ge 1$,
$\alpha_{\min} \in (0,1]$).

\end{document}
